\section{Introduction}
\label{sec:mmc-introduction}

Accurate memory measurement in Python applications is critical for large-scale scientific workloads, especially in \ac{HPC} contexts where resources must be requested in advance.
Python’s complex memory model, which mixes reference counting, dynamic allocations, and garbage collection, often obscures true memory usage.
Overestimates waste resources, whereas underestimates risk \ac{OOM} errors.

Tools like \texttt{top}\cite{top_linux_manpage}, \texttt{tracemalloc}\cite{tracemalloc}, and container-based statistics each provide a different view of memory consumption.
Containerization (\eg, Docker~\cite{docker}) further complicates accounting due to caching and shared kernel resources.
Consequently, obtaining definitive or comprehensive memory metrics in large-scale or containerized settings can be difficult.

This chapter lays out a structured approach to measuring Python memory accurately.
Section~\ref{sec:mmc-common-challenges-and-pitfalls} summarizes frequent pitfalls, such as memory fragmentation, garbage-collection spikes, and container-level distortions.
Section~\ref{sec:mmc-best-practices} consolidates best practices for environment isolation, cross-validation, and sampling strategies.
Section~\ref{sec:mmc-traceq} introduces the \texttt{traceq} library, designed to allow a memory-footprint free profiling for both Python-internal (\eg, \texttt{tracemalloc}) and system-level metrics (\eg, \texttt{psutil}).
Section~\ref{sec:mmc-recommended-strategy} walks through a step-by-step recommended measurement workflow, culminating in an example experiment that benchmarks overhead and accuracy.
Section~\ref{sec:mmc-conclusion} concludes with the main findings and outlines how these measurements underpin later chapters (notably Chapter~\ref{ch:predicting-memory-consumption-from-input-shapes} and Chapter~\ref{ch:improving-data-parallelism-using-memory-aware-chunking}).

Achieving robust, consistent, and reproducible memory metrics requires more than a single tool invocation.
It calls for isolation strategies, methodical sampling, and cross-validation.
The following sections provide those details.